\section{Einheitsvektoren}
\label{sec:vector_algebra_unit_vectors}

\begin{DefinitionBox}{Einheitsvektor}{vector_algebra_unit_vector}
  Ein \emph{Einheitsvektor} $\vec e$ \EnglishTerm{unit vector} ist ein Vektor der Länge $1$. 
  \[
    |\vec e|=1.
  \]
\end{DefinitionBox}

Häufig möchten wir zu einem Vektor den zugehörigen Einheitsvektor konstruieren, d.h.
ein Vektor der die gleiche Richtung hat, aber nur noch die Länge 1.
\begin{TheoremBox}{Normierung eines Vektors}{vector_algebra_vector_normalization}
  Durch Multiplikation mit dem Skalar $1/|\vec a|$ normieren wir einen Vektor
  so, dass seine Länge $1$ beträgt:
  \[
    \vec e_a
    =\frac1{|\vec a|}\,\vec a
    =\frac{\vec a}{|\vec a|}
    \qquad (\vec a\ne\vec 0).
  \]
  Der resultierende Einheitsvektor $\vec e_a$ hat die Richtung von $\vec a$
  und die Länge $|\vec e_a|=1$.
\end{TheoremBox}

Mit den Rechenregeln für den Betrag überprüfen wir, dass die
vorgegebene Rechenvorschrift tatsächlich einen Einheitsvektor liefert:
\begin{ClozeArea}[26]{3}
  \[
    |\vec e_a|
    =\left|\frac1{|\vec a|}\,\vec a\right|
    =\left|\frac1{|\vec a|}\right|\,|\vec a|
    =\frac1{|\vec a|}\,|\vec a|
    =1.
  \]
\end{ClozeArea}

\begin{ExampleBox}{Normierung eines Vektors}{vector_algebra_normalize_vector}
  \begin{enumerate}[label=\alph*),leftmargin=1.7em,itemsep=0.2em]
    \item
      \begin{minipage}[t]{0.66\linewidth}
        Für $\vec a=\begin{pmatrix}1\\1\end{pmatrix}$ gilt
        \[
          {\color{THMBlue}\vec e_a=}
          \begin{ClozeContinuation}[20]{3}
            {\color{THMBlue}
              \frac{\vec a}{|\vec a|}
              =\frac1{\sqrt{1^2+1^2}}\begin{pmatrix}1\\1\end{pmatrix}
              =\frac1{\sqrt2}\begin{pmatrix}1\\1\end{pmatrix}}
          \end{ClozeContinuation}.
        \]
      \end{minipage}\hfill
      \begin{minipage}[t]{0.32\linewidth}
        \vspace{0pt}
        \centering
        \CachedTikzFigure{content/resources/chapters/vector_algebra/tikz/vector_normalization_example_2d.tex}
      \end{minipage}\par

    \item Für $\vec b=\begin{pmatrix}1\\2\\2\end{pmatrix}\in\R^3$ gilt
      \[
        \vec e_b
        =\begin{ClozeContinuation}[24]{3}
          \frac{\vec b}{|\vec b|}
          =\frac1{\sqrt{1^2+2^2+2^2}}\begin{pmatrix}1\\2\\2\end{pmatrix}
          =\frac13\begin{pmatrix}1\\2\\2\end{pmatrix}
        \end{ClozeContinuation}.
      \]
  \end{enumerate}
\end{ExampleBox}

\begin{TheoremBox}{Zerlegung in Länge und Richtung}{vector_algebra_magnitude_direction_decomposition}
  Jeder Vektor $\vec a\ne\vec 0$ lässt sich als Produkt aus seinem Betrag und
  dem zugehörigen Einheitsvektor darstellen:
  \begin{center}
    \CachedTikzFigure{content/resources/chapters/vector_algebra/tikz/vector_magnitude_direction_decomposition.tex}
  \end{center}
  Damit werden die beiden Eigenschaften eines Vektors sichtbar aufgetrennt:
  \begin{itemize}[leftmargin=1.7em,itemsep=0.35em]
    \item Die \Cloze{2.2cm}{Länge} wird durch den Betrag $|\vec a|$ des Vektors beschrieben.
    \item Die \Cloze{2.4cm}{Richtung} wird durch den Einheitsvektor $\vec e_a$ beschrieben.
  \end{itemize}
\end{TheoremBox}

\clearpage
\subsection{Kanonische Einheitsvektoren}
\label{subsec:vector_algebra_canonical_unit_vectors}

Einheitsvektoren können in jede beliebige Richtung zeigen. Besonders
naheliegend sind Einheitsvektoren, die entlang der Koordinatenachsen zeigen.

\begin{DefinitionBox}{Kanonische Einheitsvektoren}{vector_algebra_canonical_unit_vectors}
  Die \emph{kanonischen Einheitsvektoren}
  \EnglishTerm{canonical unit vectors} sind die Einheitsvektoren (Länge 1) entlang der
  positiven Koordinatenachsen:
  \begin{itemize}[leftmargin=1.5em,itemsep=0.9em]
    \item Im $\R^2$:
      \par\smallskip
      \begin{minipage}[c]{0.47\linewidth}
        \[
          \begin{gathered}
            \vec e_1=\begin{pmatrix}1\\0\end{pmatrix},
            \qquad
            \vec e_2=\begin{pmatrix}0\\1\end{pmatrix},\\[0.5em]
            \text{paarweise orthogonal}:
            \vec e_1\perp\vec e_2.
          \end{gathered}
        \]
      \end{minipage}\hfill
      \begin{minipage}[c]{0.47\linewidth}
        \centering
        \CachedTikzFigure{content/resources/chapters/vector_algebra/tikz/canonical_unit_vectors_2d.tex}
      \end{minipage}

    \item Im $\R^3$:
      \par\smallskip
      \begin{minipage}[c]{0.47\linewidth}
        \[
          \begin{gathered}
            \vec e_1=\begin{pmatrix}1\\0\\0\end{pmatrix},
            \quad
            \vec e_2=\begin{pmatrix}0\\1\\0\end{pmatrix},
            \quad
            \vec e_3=\begin{pmatrix}0\\0\\1\end{pmatrix},\\[0.5em]
            \text{paarweise orthogonal:}\\[0.25em]
            \vec e_1\perp\vec e_2,
            \quad
            \vec e_1\perp\vec e_3,
            \quad
            \vec e_2\perp\vec e_3.
          \end{gathered}
        \]
      \end{minipage}\hfill
      \begin{minipage}[c]{0.47\linewidth}
        \centering
        \CachedTikzFigure{content/resources/chapters/vector_algebra/tikz/canonical_unit_vectors_3d.tex}
      \end{minipage}
  \end{itemize}
\end{DefinitionBox}

Das Fremdwort \emph{kanonisch} wird in vielen Fachgebieten für eine besonders
naheliegende oder ausgezeichnete Standardwahl verwendet.

Vektoren werden daher nicht nur als Koordinatenspalten, sondern auch mithilfe
kanonischer Einheitsvektoren dargestellt. Diese insbesondere in der
Technischen Mechanik wichtige Darstellungsform macht die Anteile in den
einzelnen Koordinatenrichtungen unmittelbar sichtbar.

\begin{ExampleBox}{Darstellung mit kanonischen Einheitsvektoren}{vector_algebra_canonical_unit_vector_decomposition}
  \begin{enumerate}[label=\alph*),leftmargin=1.7em,itemsep=1.3em]
    \item
        $\displaystyle
        \begin{pmatrix}2\\3\end{pmatrix}
        =\begin{ClozeContinuation}[28]{2}
          \begin{pmatrix}2\\0\end{pmatrix} + \begin{pmatrix}0\\3\end{pmatrix}
          = 2 \begin{pmatrix}1\\0\end{pmatrix} +3 \begin{pmatrix}0\\1\end{pmatrix}
          = 2\vec e_1+3\vec e_2
        \end{ClozeContinuation}.$

    \item
      $\displaystyle
        \begin{pmatrix}1\\2\\3\end{pmatrix}
        =\begin{ClozeContinuation}[28]{3}
          \begin{pmatrix}1\\0\\0\end{pmatrix}
          +\begin{pmatrix}0\\2\\0\end{pmatrix}
          +\begin{pmatrix}0\\0\\3\end{pmatrix}
          =1\begin{pmatrix}1\\0\\0\end{pmatrix}
          +2\begin{pmatrix}0\\1\\0\end{pmatrix}
          +3\begin{pmatrix}0\\0\\1\end{pmatrix}
          =1\vec e_1+2\vec e_2+3\vec e_3
        \end{ClozeContinuation}.$
  \end{enumerate}
  \textbf{Erkenntnis:} Die Gewichtungsfaktoren der kanonischen Einheitsvektoren $\vec{e}_i$
  entsprechen den Komponenten des Vektors.
\end{ExampleBox}

\begin{KeyPointBox}{Zerlegung in kanonische Einheitsvektoren}{vector_algebra_canonical_unit_vector_representation}
  Jeder Vektor ist die gewichtete Summe der kanonischen Einheitsvektoren:
  \[
    \begin{array}{@{}r@{:\qquad}c@{\quad=\quad}l@{}}
      \R^2
        & \begin{pmatrix}a_1\\a_2\end{pmatrix}
          \vphantom{\begin{pmatrix}a_1\\a_2\\a_3\end{pmatrix}}
        & \Cloze{10cm}{a_1\vec e_1+a_2\vec e_2},\\[1.4em]
      \R^3
        & \begin{pmatrix}a_1\\a_2\\a_3\end{pmatrix}
        & \Cloze{10cm}{a_1\vec e_1+a_2\vec e_2+a_3\vec e_3},\\[1.4em]
      \R^n
        & \begin{pmatrix}a_1\\a_2\\\vdots\\a_n\end{pmatrix}
        & \Cloze{10cm}{a_1\vec e_1+a_2\vec e_2+\dots+a_n\vec e_n}.
    \end{array}
  \]
\end{KeyPointBox}

